\clearpage
\section{ТЕОРЕТИЧЕСКИЕ ОСНОВЫ РАЗРАБОТКИ CRM-ОРИЕНТИРОВАННОГО ПРИЛОЖЕНИЯ}

\subsection{Предметная область и функциональные требования к CRM-системам}

Современная CRM-система в прикладном понимании представляет собой информационную платформу, обеспечивающую централизованный учет клиентов, историю взаимодействия, планирование активностей и контроль результативности операций. Если рассматривать учебный проект, то минимально достаточная модель должна включать сущности клиента, задачи, статусы выполнения, источники лидов и пользовательскую аналитику. Без этих компонентов невозможно обеспечить трассируемость бизнес-действий и сформировать корректную отчетную базу \cite{crmBook,projectManagementBook}.

Классическая задача CRM в цифровой среде состоит в том, чтобы перевести разрозненные процессы коммуникации в формализованные сценарии: от первичного контакта до конверсии и постпродажного сопровождения. На уровне данных это приводит к необходимости поддерживать консистентные связи между сущностями. Например, задача должна быть либо автономной, либо привязанной к конкретному клиенту, а изменение статуса клиента должно корректно отражаться в аналитических представлениях. В противном случае отчетность становится недостоверной, а управленческие решения теряют обоснованность \cite{dataDesign2022,isDesign2023}.

Для курсового проекта важным является ограничение функционального объема. С одной стороны, работа должна быть достаточно содержательной, чтобы продемонстрировать инженерные компетенции. С другой стороны, она не должна превращаться в попытку реализации корпоративной ERP-системы. Поэтому в рамках \textit{Gamified Task CRM} функциональные требования сгруппированы вокруг четырех опорных блоков:
\begin{enumerate}
  \item Пользовательский контур: авторизация, переходы между основными разделами приложения, обратная связь по ошибкам.
  \item Операционный контур: управление задачами и клиентами, поддержка CRUD-операций, назначение приоритетов, учет сроков.
  \item Аналитический контур: представление агрегированной информации по активности и прогрессу пользователя.
  \item Эксплуатационный контур: режимы выполнения (demo и backend-first), воспроизводимые проверки качества и деплой web-версии.
\end{enumerate}

Такое деление соответствует принципу функциональной декомпозиции, когда каждую подсистему можно развивать независимо при сохранении общих интерфейсных контрактов. Применение этого подхода в учебных проектах повышает управляемость кода и упрощает тестирование, поскольку сценарии проверки привязываются к четко выделенным зонам ответственности \cite{cleanArchitecture,patternsBook}.

Отдельно необходимо учитывать нефункциональные требования, которые часто игнорируются на ранних этапах обучения, но критически важны при защите проекта. К ключевым нефункциональным требованиям относятся:
\begin{itemize}
  \item надежность операций записи и чтения данных;
  \item предсказуемость пользовательских сценариев при ошибках сети и backend;
  \item масштабируемость структуры кода при добавлении новых модулей;
  \item безопасность хранения и передачи конфигурационных параметров;
  \item прозрачность контроля качества через автоматизированные проверки.
\end{itemize}

С методической точки зрения указанные требования создают мост между теорией системного анализа и практикой программной инженерии. Именно их учет в курсовом проекте позволяет перейти от формального описания интерфейсов к аргументированному выбору архитектуры.

\subsection{Архитектурные подходы и обоснование выбранного стека}

В проектировании прикладных информационных систем можно выделить два полярных подхода: монолитный сценарий с жесткой связностью компонентов и модульный сценарий с разделением ответственности. Для учебного проекта второго курса/третьего курса (в зависимости от учебного плана) предпочтителен второй вариант, поскольку он лучше демонстрирует инженерную зрелость и позволяет управлять сложностью разработки \cite{sommerville,designPatterns}.

С учетом поставленных задач в \textit{Gamified Task CRM} используется кроссплатформенный UI-фреймворк Flutter, язык Dart, механизм управления состоянием Riverpod, декларативная маршрутизация GoRouter и backend-платформа Supabase. Ниже приведено сравнительное обоснование такого выбора.

\begin{table}[h!]
\caption{Сравнение ключевых технологических решений проекта}
\centering
\begin{tabular}{|p{4cm}|p{4.5cm}|p{5.7cm}|}
\hline
\textbf{Компонент} & \textbf{Рассмотренные варианты} & \textbf{Причина выбора в проекте} \\
\hline
UI-фреймворк & Flutter, React + TypeScript, Vue & Единая кодовая база, быстрый UI-прототипинг, зрелая экосистема для web и мобильного расширения \cite{flutterDocs} \\
\hline
Управление состоянием & setState, BLoC, Riverpod & Декларативность, тестируемость, удобная композиция провайдеров, меньшая связность \cite{riverpodDocs} \\
\hline
Маршрутизация & Navigator 1.0, auto\_route, GoRouter & Прозрачные redirect-правила, удобная интеграция с авторизацией и deep-link сценариями \cite{goRouterDocs} \\
\hline
Backend & Firebase, Supabase, custom API & PostgreSQL-модель данных, RLS-политики, SQL-миграции и удобный локальный/облачный контур \cite{supabaseDocs,postgresDocs} \\
\hline
\end{tabular}
\label{tab:tech-comparison}
\end{table}

Данные таблицы \ref{tab:tech-comparison} показывают, что выбор стека основан не на популярности инструментов, а на соответствии требованиям проекта: управляемость состояния, прозрачная маршрутизация, поддержка реляционной модели и доступность практик безопасности. Для курсовой работы такой подход является корректным, поскольку демонстрирует умение сопоставлять техническое решение с предметной задачей.

С архитектурной точки зрения применим принцип \textit{feature-first organization}, когда код группируется по функциональным модулям (dashboard, tasks, crm, reports, profile, ai), а не по техническим слоям в изоляции. Это решение улучшает локальность изменений: добавление нового поведения в одном модуле требует минимального вмешательства в соседние. При этом общие инфраструктурные элементы (конфигурация, база, авторизация) остаются централизованными и переиспользуемыми \cite{repoReadme,patternsBook}.

Для формального описания эксплуатационной результативности можно использовать показатель покрытия ключевых сценариев:
\begin{equation}
Q = \frac{N_{pass}}{N_{total}}.
\label{eq:quality}
\end{equation}

где $N_{pass}$ --- количество успешно пройденных проверок, а $N_{total}$ --- общее количество запущенных проверок в контрольном контуре. Показатель \eqref{eq:quality} не исчерпывает понятие качества, но служит удобной метрикой для оперативного мониторинга технической стабильности при развитии проекта \cite{testingBook}.

\subsection{Теоретические основы проектирования данных и клиентских сценариев}

Проектирование модели данных является центральной задачей для любой CRM-системы, поскольку именно данные связывают пользовательские действия, бизнес-логику и аналитические отчеты. В рамках рассматриваемого проекта ключевая логическая связь задается через отношение \textit{tasks.client\_id $\rightarrow$ crm\_clients.id}. Такая связь допускает как привязанные, так и автономные задачи, что отражает реальный процесс работы: часть задач связана с конкретным клиентом, часть относится к внутренним активностям команды \cite{repoReadme,migrationClientId}.

При построении реляционной модели для учебного проекта целесообразно придерживаться следующих правил:
\begin{enumerate}
  \item Каждая сущность должна иметь устойчивый первичный идентификатор.
  \item Внешние связи должны быть явными и проверяемыми на уровне схемы данных.
  \item Поля с неопределенным значением должны иметь корректную nullable-семантику.
  \item Дата-время в операционных таблицах должна сопровождаться полями создания и обновления.
\end{enumerate}

Указанные правила соответствуют базовым принципам проектирования информационных систем и существенно упрощают последующую интеграцию прикладной логики с хранилищем. Например, наличие полей \textit{created\_at} и \textit{updated\_at} позволяет организовать сортировку, аудит изменений и формирование периодических отчетов без дополнительной денормализации \cite{dbBook2021,postgresDocs}.

С точки зрения пользовательских сценариев важна не только структурная корректность данных, но и целостность переходов между экранами. Типовой процесс CRM-пользователя можно представить в виде цепочки: вход в систему, создание или выбор клиента, постановка задачи, контроль исполнения, просмотр сводной информации. Наличие разрывов в этой цепочке снижает практическую ценность приложения, поскольку пользователь вынужден выполнять внешние действия вне системы.

В теории UX-проектирования это описывается как требование \textit{continuity of workflow}: пользователь должен завершить ключевой сценарий без неоправданных переключений контекста \cite{uxBook2023}. В проекте \textit{Gamified Task CRM} данное требование поддерживается за счет единой оболочки с нижней навигацией и согласованных форм добавления сущностей. С инженерной точки зрения это также уменьшает вероятность ошибок, связанных с дублированием логики в независимых экранах.

На рисунке \ref{fig:architecture} показана обобщенная схема модульного взаимодействия в системе.

\begin{figure}[h!]
\centering
\fbox{%
\begin{minipage}{0.9\textwidth}
\centering
UI-модули (Dashboard / Tasks / CRM / Reports / Profile / AI) \\
$\downarrow$ \\
Провайдеры состояния (Riverpod Notifiers) \\
$\downarrow$ \\
Сервисы данных (Supabase API, Auth, локальный demo-контур) \\
$\downarrow$ \\
PostgreSQL-таблицы и политики доступа
\end{minipage}}
\caption{Обобщенная архитектурная схема проекта}
\label{fig:architecture}
\end{figure}

Как показано на рисунке \ref{fig:architecture}, модульность достигается не только разбиением интерфейса, но и четким разделением уровней ответственности: UI, состояние, сервисы, хранилище. Такой подход упрощает сопровождение и снижает риск каскадных дефектов при изменении одного из компонентов.

\subsection{Качество программного продукта: тестирование, надежность, безопасность}

Качество программного продукта в инженерной практике интерпретируется как совокупность нескольких свойств: функциональная корректность, устойчивость к ошибкам, эксплуатационная предсказуемость и безопасность. Для курсового проекта важно показать, что эти свойства учитываются не декларативно, а через конкретные процедуры проверки \cite{testingBook,owaspM10}.

В рассматриваемой системе используется многоуровневый подход к верификации:
\begin{enumerate}
  \item \textbf{Статический анализ} --- выявление потенциальных ошибок и нарушений стиля до выполнения приложения.
  \item \textbf{Unit/Widget-тесты} --- проверка отдельных компонентов и UI-поведения в изолированной среде.
  \item \textbf{Integration smoke} --- контроль целостности основного бизнес-потока после изменений.
  \item \textbf{Web E2E} --- сквозная проверка критичного сценария в среде, максимально близкой к эксплуатационной.
\end{enumerate}

Теоретически такой набор тестов соответствует принципу \textit{test pyramid}: широкая база быстрых проверок на нижних уровнях и ограниченное число более дорогих сквозных сценариев на верхнем уровне \cite{testingBook}. Для учебного проекта это особенно важно, так как позволяет соблюдать баланс между глубиной проверки и временем выполнения.

Отдельную роль играет обработка исключительных ситуаций. Даже корректно спроектированная архитектура теряет практическую ценность, если пользователь получает неинформативные сообщения об ошибках. В рамках проекта используется нормализация технических backend-ошибок в человеко-понятные сообщения (например, отсутствие таблицы, проблема RLS, истекшая сессия). Теоретически это относится к паттерну \textit{error translation boundary}, когда низкоуровневые ошибки преобразуются на границе слоя сервиса в формат, пригодный для пользовательского интерфейса \cite{cleanArchitecture,patternsBook}.

С позиции информационной безопасности в курсовом проекте применяются базовые, но принципиальные практики:
\begin{itemize}
  \item вынесение секретов из кода в параметры окружения;
  \item разделение demo и production режимов;
  \item ограничение доступа к данным через политики уровня строк;
  \item контроль аутентификации перед операциями записи.
\end{itemize}

Эти меры соответствуют рекомендациям OWASP по предотвращению типичных ошибок конфигурации и контроля доступа в web-приложениях \cite{owaspM4,owaspM1}. Для учебной разработки их достаточно, чтобы продемонстрировать осознанный подход к защищенной реализации.

\subsection{Теоретические выводы по разделу}

Проведенный анализ показывает, что разработка CRM-ориентированного приложения требует системного сочетания трех групп решений: предметных (что делает система), архитектурных (как организован код) и эксплуатационных (как подтверждается корректность и устойчивость). Игнорирование любой из этих групп приводит либо к функционально неполному продукту, либо к трудно сопровождаемой реализации.

Для темы курсового проекта обоснован выбор Flutter-стека с Riverpod, GoRouter и Supabase, поскольку он обеспечивает:
\begin{itemize}
  \item модульную организацию кода и масштабируемость;
  \item прозрачные пользовательские сценарии;
  \item реляционную модель данных с контролируемыми связями;
  \item возможность поэтапной верификации качества.
\end{itemize}

Теоретическая часть также подтверждает, что качество итогового решения определяется не только полнотой функционала, но и зрелостью инженерных практик: корректной конфигурацией окружения, продуманной обработкой ошибок, повторяемыми тестовыми сценариями и соблюдением минимальных требований безопасности.

Сформированные в разделе положения являются методической базой для практической главы, где рассматривается реализация \textit{Gamified Task CRM}, описание архитектуры в коде, структура данных, сценарии использования и результаты проверки работоспособности.
